\section*{Abstract}
		Diese Arbeit präsentiert eine umfassende Analyse der Temperatureinheiten in natürlichen Einheiten ($\hbar = c = k_B = 1$) im Rahmen der T0-Theorie. Im statischen $\xi$-Universum wird keine expandierende Raumzeit benötigt. Die Ableitungen verwenden die universelle Konstante $\xi = \frac{4}{3} \times 10^{-4}$ als einzigen Parameter (plus die im Korpus genannten ganzzahligen Setzungen; SI-Werte dienen der Umrechnung) und respektieren die fundamentale Zeit-Energie-Dualität. Das Dokument beinhaltet CMB-Berechnungen im Rahmen der T0-Theorie, behandelt fundamentale Fragen zu Rotverschiebungsmechanismen, primordialen Störungen und der Auflösung kosmologischer Spannungen. Die Theorie deutet die CMB ohne Inflation und ohne eine Entkopplung bei irgendeinem kosmologischen $z$ -- durch stationäre Thermalisierung des $\xi$-Feldes in einem unendlich alten Universum --, leitet primordiale Störungen aus T-Feld-Quantenfluktuationen ab und behandelt die Hubble-Spannung (siehe Dokument~026 für die an $H_0$ kalibrierte Beschreibung mit $H_0 \approx 66{,}2\,$km/s/Mpc; der kosmische Sektor ist kalibriert, nicht aus $\xi$ hergeleitet).

	\medskip\noindent\textbf{Statusmarker.} Aussagen mit \textbf{[S]} sind \emph{spekulativ}: ein Hinweis oder eine Vermutung, die im Korpus noch nicht hergeleitet oder numerisch geprüft ist. Der heutige Korpus klammert den kosmischen Sektor ($H_0$, $\Lambda$, Dunkle Energie, Ursprung der CMB) bewusst aus, weil auch das Standardmodell ihn nicht herleitet (P39); die kosmologischen Deutungen dieses Dokuments (statisches, ewiges Universum, Ursprung der CMB, Vergleich mit $\Lambda$CDM) sind in diesem Sinn zu lesen.

	
	
	\section{Einführung: T0-Theorie in natürlichen Einheiten}
	
	\subsection{Natürliche Einheiten als Grundlage}
	
	\begin{important}
		Diese gesamte Arbeit verwendet ausschließ{}lich natürliche Einheiten mit $\hbar = c = k_B = 1$. Alle Größ{}en haben Energiedimensionen: $[L] = [T] = [E^{-1}]$, $[M] = [T_{\text{temp}}] = [E]$.
	\end{important}
	
	Das System der natürlichen Einheiten stellt eine fundamentale Vereinfachung der Physik dar, indem die universellen Konstanten $\hbar$ (reduzierte Planck-Konstante), $c$ (Lichtgeschwindigkeit) und $k_B$ (Boltzmann-Konstante) auf den Wert 1 gesetzt werden. Diese Wahl vereinfacht die Formeln und macht die Beziehungen zwischen den Größ{}en direkt sichtbar.
	
	In diesem System reduziert sich die gesamte Physik auf eine einzige fundamentale Dimension - Energie. Alle anderen physikalischen Größ{}en werden als Potenzen der Energie ausgedrückt:
	\begin{align}
		\text{Länge:} \quad [L] &= [E^{-1}] \quad \text{(Energie}^{-1}\text{)} \\
		\text{Zeit:} \quad [T] &= [E^{-1}] \quad \text{(Energie}^{-1}\text{)} \\
		\text{Masse:} \quad [M] &= [E] \quad \text{(Energie)} \\
		\text{Temperatur:} \quad [T_{\text{temp}}] &= [E] \quad \text{(Energie)}
	\end{align}
	
	Diese dimensionale Reduktion macht Beziehungen zwischen Größ{}en transparent. In natürlichen Einheiten wird beispielsweise Einsteins berühmte Formel $E = mc^2$ zur trivialen Aussage $E = m$, da sowohl Energie als auch Masse dieselbe Dimension haben.
	
	\textbf{Einheitenumrechnung (zur Referenz):}
	Für Leser, die mit SI-Einheiten vertraut sind, gelten folgende Umrechnungsfaktoren:
	\begin{itemize}
		\item $\hbar = 1{,}055 \times 10^{-34}$ J$\cdot$s $\rightarrow 1$ (nat. Einheiten)
		\item $c = 2{,}998 \times 10^8$ m/s $\rightarrow 1$ (nat. Einheiten)  
		\item $k_B = 1{,}381 \times 10^{-23}$ J/K $\rightarrow 1$ (nat. Einheiten)
	\end{itemize}
	
	\subsection{Die universelle $\xi$-Konstante}
	
	\begin{important}[Kernaussage]
		Die T0-Theorie verwendet eine einzige geometrische Konstante $\xi = \frac{4}{3} \times 10^{-4}$ als Grundparameter; aus ihr werden Größ{}en von den Quarks bis zu kosmischen Strukturen abgeleitet. Kosmologisch beschreibt dieses Dokument einen statischen, ewig existierenden Kosmos ohne Urknall \textbf{[S]}. Der Faktor $\frac{4}{3}$ stammt aus dem fundamentalen geometrischen Verhältnis zwischen Kugelvolumen und Tetraedervolumen im dreidimensionalen Raum.
	\end{important}
	
	Das Herz der T0-Theorie bildet eine universelle dimensionslose Konstante, die wir mit dem griechischen Buchstaben $\xi$ (Xi) bezeichnen. Diese Konstante wurde ursprünglich rein geometrisch aus den fundamentalen T0-Feldgleichungen abgeleitet, wie in der etablierten T0-Theorie \cite{T0Theory} gezeigt.
	
	Die fundamentale T0-Theorie basiert auf der universellen dimensionslosen Konstante:
	\begin{equation}
		\xi = \frac{4}{3} \times 10^{-4} \quad \text{(dimensionslos, exakter geometrischer Wert)}
	\end{equation}
	
	\textbf{Geometrische Ableitung aus T0-Feldgleichungen:} Der Wert von $\xi$ folgt direkt aus der geometrischen Struktur der T0-Feldgleichungen des universellen Energiefeldes $E_{\text{field}}(x,t)$. Die fundamentale T0-Gleichung $\square E_{\text{field}} = 0$ in Verbindung mit dreidimensionaler Raumgeometrie führt zu:
\begin{itemize}
	\item Der geometrische Faktor $\frac{4}{3}$ aus der dreidimensionalen Raumgeometrie
	\item Das Skalenverhältnis $10^{-4}$ aus der fraktalen Dimension
	\item Für die vollständige Herleitung siehe 041\_parameterherleitung\_De.pdf 
\end{itemize}
	
	\textbf{Vergleich mit dem Experiment:} Die Abweichung des anomalen magnetischen Moments des Myons vom Standardmodell ist mit dem Stand 2025 nicht mehr signifikant; sie dient daher nicht als Beleg für $\xi$ (Dok.~382). Die $g-2$-Linie des Korpus beruht auf Verhältnissen der Leptonen; ihre Testaussage ist die an gemessenen Größ{}en verankerte Tau-Vorhersage $a_\tau \approx 1{,}281\times10^{-3}$ (Dok.~382).
	
	Diese Konstante bestimmt in der T0-Theorie eine Vielzahl physikalischer Phänomene:
	\begin{itemize}
		\item \textbf{Teilchenphysik}: Alle Elementarteilchenmassen ergeben sich aus geometrischen Quantenzahlen $(n,l,j,r,p)$ skaliert mit $\xi$
		\item \textbf{Feldtheorie}: Charakteristische Energieskalen aller Wechselwirkungen folgen aus $\xi$-Felddynamik
		\item \textbf{Gravitation}: Die Gravitationskonstante in natürlichen Einheiten $G_{\text{nat}} = 2{,}61 \times 10^{-70}$ ist eine direkte Funktion von $\xi$
		\item \textbf{Kosmologie} \textbf{[S]}: Thermodynamisches Gleichgewicht im statischen, unendlich alten Universum wird durch $\xi$-Feldzyklen aufrechterhalten
	\end{itemize}
	
	\textbf{Symbolerklärung:}
	\begin{itemize}
		\item $\xi$ (Xi): Universelle dimensionslose Konstante der T0-Theorie
		\item $E_\xi$: Charakteristische Energieskala, definiert als $E_\xi = 1/\xi$
		\item $T_\xi$: Charakteristische Temperatur, gleich $E_\xi$ in natürlichen Einheiten
		\item $L_\xi$: Charakteristische Längenskala des $\xi$-Feldes
		\item $G_{\text{nat}}$: Gravitationskonstante in natürlichen Einheiten
		\item $\alpha_{\text{EM}}$: Elektromagnetische Kopplung (= 1 in natürlichen Einheiten per Definition)
		\item $\beta$: Dimensionsloser Parameter $\beta = r_0/r = 2GE/r$
		\item $\omega$: Photonenenergie (Dimension $[E]$ in natürlichen Einheiten)
	\end{itemize}
	
	\textbf{Kopplungskonstanten in natürlichen Einheiten:}
	\begin{align}
		\alpha_{\text{EM}} &= 1 \quad \text{(per Definition in natürlichen Einheiten)} \\
		\alpha_G &= \xi^2 = \left(\frac{4}{3} \times 10^{-4}\right)^2 = 1{,}78 \times 10^{-8} \\
		\alpha_W &= \xi^{1/2} = \left(\frac{4}{3} \times 10^{-4}\right)^{1/2} = 1{,}15 \times 10^{-2} \\
		\alpha_S &= \xi^{-1/3} = \left(\frac{4}{3} \times 10^{-4}\right)^{-1/3} = 19{,}57
	\end{align}
	Zum Vergleich: gemessen ist $\alpha_s(M_Z) \approx 0{,}118$.
	
	\textbf{Wichtige Klarstellung zu Einheiten:}
	In diesem gesamten Dokument arbeiten wir ausschließ{}lich in natürlichen Einheiten mit $\hbar = c = k_B = 1$. Das bedeutet:
	\begin{itemize}
		\item Die elektromagnetische Kopplungskonstante ist $\alpha_{\text{EM}} = 1$ per Definition (nicht 1/137 wie in SI-Einheiten)
		\item Alle anderen Kopplungskonstanten werden relativ zu $\alpha_{\text{EM}} = 1$ ausgedrückt
		\item Energie, Masse und Temperatur haben dieselbe Dimension
		\item Länge und Zeit haben die Dimension Energie$^{-1}$
	\end{itemize}
	
	\textbf{Dimensionale Konsistenz:} Da $\xi$ rein dimensionslos ist, hat es denselben Wert in allen Einheitensystemen. Es charakterisiert in der T0-Theorie die Geometrie des Raum-Zeit-Kontinuums und wird als Naturkonstante behandelt, vergleichbar mit der Feinstrukturkonstante.
	
	\subsection{Zeit-Energie-Dualität und statisches Universum}
	
	\begin{important}
		\textbf{[S]} Dieses Dokument deutet das Universum als ewig existierend, ohne zeitlichen Anfang. Die Begründung, dass es keinen singulären Anfang gibt, folgt aus der Zeit-Energie-Dualität $T\cdot E = 1$ mit der unteren Zeitschranke $T_0 = \xi\,t_P$ und dem randlosen Torus (Dok.~306). Die kosmologische Deutung bleibt eine Deutung, kein Beweis; der kosmische Sektor ist im heutigen Korpus bewusst ausgeklammert (P39).
	\end{important}
	
	Aus der Zeit-Masse-Dualität $T\cdot m = 1$ (Dok.~003/010) folgt in natürlichen Einheiten die Zeit-Energie-Dualität
	\begin{equation}
		T \cdot E = 1 .
	\end{equation}
	
	Damit gilt $E\to\infty \Leftrightarrow T\to 0$. Die Zeit ist in der T0-Theorie jedoch nach unten beschränkt durch $T_0 = \xi\,t_P$, die Zeit-Kleidung der minimalen Länge $L_0 = \xi\,\ell_P$; entsprechend ist die Energie nach oben beschränkt durch $E_\text{max} = E_P/\xi$, mit exakter Schließung $T_0\,E_\text{max} = 1$. Eine unendliche Energiedichte ist damit strukturell ausgeschlossen. Ergänzend ist der Torus kompakt und randlos, $\partial T^4 = \emptyset$: Die Frage nach einem Anfang bei $t=0$ setzt einen Rand voraus, den ein kompakter Torus nicht besitzt (Dok.~306).
	
	Heisenbergs Energie-Zeit-Unschärferelation $\Delta E \times \Delta t \geq 1/2$ trägt diesen Schluss nicht: Ein endliches $\Delta t$ liefert eine \emph{endliche} Schranke für $\Delta E$, nicht $\Delta E\to\infty$, und das $\Delta t$ der Relation ist die Änderungszeit einer Observablen, kein kosmisches Alter (Dok.~306).
	
	\textbf{Folgerung in diesem Dokument [S]:} Es gibt keinen singulären Anfang; dieses Dokument deutet das Universum daher als ewig existierend. Dies führt zum statischen T0-Universum, das folgende Eigenschaften besitzt:
	
	Das T0-Universum ist daher:
	\begin{itemize}
		\item \textbf{Statisch}: Kein expandierender Raum - die Raumzeitmetrik ist zeitunabhängig
		\item \textbf{Ewig}: Ohne zeitlichen Anfang oder Ende
		\item \textbf{Thermodynamisch ausgeglichen}: Durch $\xi$-Feldzyklen wird ein dynamisches Gleichgewicht aufrechterhalten
		\item \textbf{Strukturell stabil}: Kontinuierliche Bildung und Erneuerung von Materie und Strukturen
	\end{itemize}
	
	\textbf{Einheitenprüfung der Zeit-Energie-Dualität:}
	\begin{align}
		[T] \times [E] &= [E^{-1}] \times [E] = [E^0] = \text{dimensionslos} \\
		[1] &= \text{dimensionslos} \quad \text{(erfüllt)}
	\end{align}
	
	\section{$\xi$-Feld und charakteristische Energieskalen}
	
	\subsection{$\xi$-Feld als universeller Energievermittler}
	
	\begin{formula}
		Die universelle Konstante $\xi = \frac{4}{3} \times 10^{-4}$ definiert die fundamentale Energieskala der T0-Theorie:
		\begin{equation}
			E_\xi = \frac{1}{\xi} = \frac{1}{\frac{4}{3} \times 10^{-4}} = \frac{3}{4} \times 10^4 = 7500
		\end{equation}
		(alle Größ{}en in natürlichen Einheiten)
	\end{formula}
	
	Das $\xi$-Feld repräsentiert das fundamentale Energiefeld des Universums, aus dem alle anderen Felder und Wechselwirkungen hervorgehen. Seine charakteristische Energieskala $E_\xi$ ergibt sich als Kehrwert der dimensionslosen Konstante $\xi$.
	
	\textbf{Einheitenprüfung für $E_\xi$:}
	\begin{align}
		[E_\xi] &= \left[\frac{1}{\xi}\right] = \frac{[E^0]}{[E^0]} = [E^0] = \text{dimensionslos}
	\end{align}
	
	In natürlichen Einheiten ist dimensionslos äquivalent zu einer Energieeinheit, da alle Größ{}en auf Energiepotenzen reduziert werden. Daher gilt $[E_\xi] = [E]$.
	
	Diese charakteristische Energie entspricht direkt einer charakteristischen Temperatur in natürlichen Einheiten, da Energie und Temperatur dieselbe Dimension haben:
	\begin{equation}
		T_\xi = E_\xi = \frac{3}{4} \times 10^4 = 7500 \quad \text{(nat. Einheiten)}
	\end{equation}
	
	\textbf{Einheitenprüfung für $T_\xi$:}
	\begin{align}
		[T_\xi] = [E_\xi] = [E] = [T_{\text{temp}}] \quad \text{(erfüllt)}
	\end{align}
	
	\textbf{Physikalische Interpretation:} Die Energieskala $E_\xi = 7500$ in natürlichen Einheiten entspricht einer extrem hohen Temperatur, die charakteristisch für die fundamentalen Prozesse des $\xi$-Feldes ist. Diese Energie liegt weit über allen bekannten Teilchenenergien.
	
	\subsection{Charakteristische $\xi$-Längenskala}
	
	Das $\xi$-Feld definiert auch eine charakteristische Längenskala:
	\begin{equation}
		L_\xi = \frac{1}{E_\xi} = \frac{1}{7500} \approx 1,33 \times 10^{-4} \quad \text{(nat. Einheiten)}
	\end{equation}
	
	Diese Längenskala spielt eine fundamentale Rolle in der geometrischen Struktur der Raumzeit und erscheint in verschiedenen physikalischen Phänomenen.
	
	\section{CMB in der T0-Theorie: Statisches $\xi$-Universum}
	
	\subsection{CMB ohne Urknall}
	
	\begin{important}[Kosmologischer Sektor]
		\textbf{[S]} Nach der obigen Argumentation schließt die Zeit-Energie-Dualität einen Urknall aus; die CMB-Hintergrundstrahlung hätte dann einen anderen Ursprung als die z=1100-Entkopplung. Diese Einordnung ist spekulativ (P39).
	\end{important}
	
	Die T0-Theorie deutet die kosmische Mikrowellen-Hintergrundstrahlung durch $\xi$-Feld-Mechanismen \textbf{[S]}:
	
	\subsubsection{1. $\xi$-Feld-Quantenfluktuationen}
	Das allgegenwärtige $\xi$-Feld erzeugt Vakuumfluktuationen mit charakteristischer Energieskala. Die exakte Abhängigkeit wird durch das gemessene Verhältnis $T_{\text{CMB}}/E_\xi \approx \xi^2$ abgeleitet.
	
	\subsubsection{2. Stationäre Thermalisierung}
	In einem unendlich alten Universum erreicht die Hintergrundstrahlung ein thermodynamisches Gleichgewicht bei der charakteristischen $\xi$-Temperatur.
	
	\begin{sibox}
		\textbf{CMB-Messungen (nur zur Referenz, in SI-Einheiten):}
		\begin{itemize}
			\item Vakuumenergiedichte: $\rho_{\text{Vakuum}} = 4,17 \times 10^{-14}$ J/m$^3$
			\item Strahlungsleistung: $j = 3,13 \times 10^{-6}$ W/m$^2$
			\item Temperatur: $T = 2,7255$ K
		\end{itemize}
	\end{sibox}
	
	\subsection{Die bereits etablierte $\xi$-Geometrie}
	
	\begin{important}
		Die T0-Theorie hatte bereits eine fundamentale Längenskala etabliert, bevor die CMB-Analyse durchgeführt wurde. Die CMB-Energiedichte wird im Folgenden mit dieser bereits existierenden $\xi$-geometrischen Struktur verglichen.
	\end{important}
	
	Aus der ursprünglichen T0-Theorie-Formulierung folgte:
	
	\textbf{Charakteristische Masse:}
	\begin{equation}
		m_{\text{char}} = \frac{\xi}{2\sqrt{G_{\text{nat}}}} \approx 4,13 \times 10^{30} \quad \text{(nat. Einheiten)}
	\end{equation}
	
	\textbf{Universelle Skalierungsregel:}
	\begin{equation}
		\text{Faktor} = 2,42 \times 10^{-31} \cdot m \quad \text{(für beliebige Masse } m \text{ in nat. Einheiten)}
	\end{equation}
	
	\textbf{Gravitationskonstante abgeleitet aus $\xi$:}
	\begin{equation}
		G_{\text{nat}} = 2,61 \times 10^{-70} \quad \text{(nat. Einheiten)}
	\end{equation}
	
	% ================== VOLLSTÄNDIGER CMB-ABSCHNITT AUS 039_Zwei-Dipole-CMB_De.pdf ==================
	
	\section{Das T0-Theorie-Rahmenwerk für CMB}
	\label{sec:t0_framework}
	
	Die T0-Theorie schlägt eine Erweiterung der Standardkosmologie durch die Einführung eines intrinsischen Zeitfeldes $\Tfield$ vor, das an alle Materie und Strahlung koppelt. Diese Theorie entstand aus der Unzufriedenheit mit der quantenmechanischen Nichtlokalität und dem Bedürfnis nach einem deterministischen Rahmenwerk, das die Kausalität bewahrt und gleichzeitig beobachtete Korrelationen erklärt.
	
	\subsection{Fundamentale Postulate}
	
	Die T0-Theorie basiert auf drei fundamentalen Postulaten:
	
	\begin{enumerate}
		\item \textbf{Zeit-Masse-Dualität}: Die fundamentale Beziehung
		\begin{equation}
			\Tfield \cdot m(x) = 1
			\label{eq:time_mass_duality}
		\end{equation}
		
		\item \textbf{Universeller Kopplungsparameter}: Ein einzelner Parameter
		\begin{equation}
			\xipar = \frac{4}{3} \times 10^{-4}, \qquad \frac{\lambda_h^2 v^2}{16\pi^3 m_h^2} \approx 1{,}30 \times 10^{-4}
			\label{eq:xi_definition}
		\end{equation}
		an der Higgs-Physik auf rund $2{,}3\,\%$ geprüft, regiert alle T-Feld-Wechselwirkungen. Mit der SM-Selbstkopplung $\lambda_h = m_h^2/(2v^2) \approx 0{,}129$ gilt $\lambda_h^2 v^2/(16\pi^3 m_h^2) = m_h^2/(64\pi^3 v^2) \approx 1{,}30 \times 10^{-4}$, rund $2{,}3\,\%$ unter $4/30000$ (PDG 2024: $m_h = 125{,}20$~GeV, $v = 246{,}22$~GeV); das ist eine Konsistenzprüfung von $\xi$, keine exakte Herleitung (Herkunft in Dok.~385). Der Faktor $\frac{4}{3}$ stammt letztendlich aus dem fundamentalen geometrischen Verhältnis zwischen Kugelvolumen und Tetraedervolumen im dreidimensionalen Raum.
		
		\item \textbf{Modifizierte Robertson-Walker-Metrik}:
		\begin{equation}
			ds^2 = -c^2dt^2[1 + 2\xipar\ln(a)] + a^2(t)[1 - 2\xipar\ln(a)]d\vec{x}^2
			\label{eq:modified_metric}
		\end{equation}
	\end{enumerate}
	
	\section{Leistungsspektren-Berechnungen}
	\label{sec:power_spectra}
	
	\subsection{Temperatur-Leistungsspektrum}
	
	Das CMB-Temperatur-Leistungsspektrum ist:
	
	\begin{equation}
		C_\ell^{TT} = \frac{2}{\pi}\int_0^\infty k^2 dk \, \mathcal{P}_\Psi(k) |\Theta_\ell(k,\eta_0)|^2 \times \left(1 + \xipar f_\ell(k)\right)
		\label{eq:cl_tt}
	\end{equation}
	
	wobei:
	\begin{equation}
		f_\ell(k) = \ln^2\left(\frac{k}{k_*}\right) - 2\ln\left(\frac{k}{k_*}\right)
	\end{equation}
	
	\subsection{E-Modus-Polarisation}
	
	\begin{equation}
		C_\ell^{EE} = \frac{2}{\pi}\int_0^\infty k^2 dk \, \mathcal{P}_\Psi(k) |E_\ell(k,\eta_0)|^2 \times \left(1 + \xipar g_\ell(k)\right)
	\end{equation}
	
	\subsection{Kreuzkorrelation}
	
	\begin{equation}
		C_\ell^{TE} = \frac{2}{\pi}\int_0^\infty k^2 dk \, \mathcal{P}_\Psi(k) \Theta_\ell(k,\eta_0) E_\ell^*(k,\eta_0) \times \left(1 + \xipar h_\ell(k)\right)
	\end{equation}
	
	\section{MCMC-Analyse und Parameter-Einschränkungen}
	\label{sec:mcmc}
	
	\subsection{Bayessche Parameter-Schätzung}
	
	Wir führen eine MCMC-Analyse durch mit:
	
	\begin{equation}
		\mathcal{L} = -\frac{1}{2}\sum_{\ell} \frac{2\ell+1}{2} f_{\text{sky}} \left[\frac{C_\ell^{\text{obs}} - C_\ell^{\text{theory}}(\theta)}{\sigma_\ell}\right]^2
	\end{equation}
	
	\subsection{Ergebnisse mit Unsicherheiten}
	
	\begin{table}[htbp]
		\centering
		\caption{T0-Parameter-Einschränkungen (68\% CL)}
		\adjustbox{max width=\textwidth}{%
\begin{tabular}{lcc}
			\toprule
			Parameter & Beste Anpassung & Unsicherheit \\
			\midrule
			$H_0$ [km/s/Mpc] & 66,2$^*$ & an $H_0$ kalibriert \\
			$\Omega_b h^2$ & 0,02237 & $\pm 0,00015$ \\
			$\Omega_c h^2$ & 0,1200 & $\pm 0,0012$ \\
			$\tau$ & 0,054 & $\pm 0,007$ \\
			$n_s$ & 0,9649 & $\pm 0,0042$ \\
			$\ln(10^{10}A_s)$ & 3,044 & $\pm 0,014$ \\
			$\xipar$ & $\frac{4}{3} \times 10^{-4}$ & (geometrische Konstante) \\
			\bottomrule
		\end{tabular}%
}
		\label{tab:parameters}
	\end{table}
{$^*$Die an $H_0$ kalibrierte T0-Beschreibung $H_0^{\text{T0}} = E_H/\hbar \approx 66{,}2\,$km/s/Mpc ($-1{,}9\%$ Abweichung von Planck) sowie der dort vorgeschlagene Auflösungsweg für die Hubble-Spannung sind in Dokument~026 (Geometrische Kosmologie) zentralisiert.}
	
	\subsection{$S_8$-Spannung}
	
	Die Clustering-Amplitude wird modifiziert:
	
	\begin{equation}
		S_8^{T0} = S_8^{\Lambda\text{CDM}} \times (1 - 2\xipar) = 0,834 \times (1 - 2 \times \frac{4}{3} \times 10^{-4}) = 0,834 \times 0,99973 = 0,8338
	\end{equation}
	
	Nach Einschätzung dieses Dokuments ist das mit schwachen Linsenmessungen verträglich \textbf{[S]}.
	
	\section{Experimentelle Vorhersagen}
	\label{sec:predictions}
	
	\subsection{Testbare Vorhersagen}
	
	Die T0-Theorie macht mehrere spezifische Vorhersagen; als kosmologische Aussagen sind sie spekulativ \textbf{[S]} (P39):
	
	\begin{enumerate}
		\item \textbf{Laufen des spektralen Index}:
		\begin{equation}
			\frac{dn_s}{d\ln k} = -2\xipar = -2 \times \frac{4}{3} \times 10^{-4} = -2,67 \times 10^{-4}
		\end{equation}
		
		\item \textbf{Tensor-zu-Skalar-Verhältnis}:
		\begin{equation}
			r = 16\xipar = 16 \times \frac{4}{3} \times 10^{-4} = 0,00213 \pm 0,0004
		\end{equation}
		
		\item \textbf{Modifizierte Silk-Dämpfung}:
		\begin{equation}
			C_\ell^{TT} \propto \exp\left[-\left(\frac{\ell}{\ell_D}\right)^2\right] \times \left(1 + \xipar \left(\frac{\ell}{3000}\right)^2\right)
		\end{equation}
	\end{enumerate}
	
	Eine wellenlängenabhängige Rotverschiebung $\Delta z \propto \ln(\lambda/\lambda_0)$ wird nicht vorhergesagt: Die fraktale Wegverlängerung ist rein geometrisch und streckt alle Wellenlängen mit demselben Faktor, die Rotverschiebung ist achromatisch (K6).
	
	\subsection{Beobachtungstests}
	
	\begin{table}[htbp]
		\centering
		\caption{T0-Vorhersagen vs Beobachtungen}
		\adjustbox{max width=\textwidth}{%
\begin{tabular}{p{3cm}p{3cm}p{3cm}p{3cm}}
			\toprule
			Beobachtbare & T0-Vorhersage & Aktuelle Grenze & Zukünftige Sensitivität \\
			\midrule
			$dn_s/d\ln k$ & $-2,67 \times 10^{-4}$ & $< 0,01$ & $10^{-4}$ (CMB-S4) \\
			$r$ & $0,00213$ & $< 0,036$ & $0,001$ (LiteBIRD) \\
			$f_{NL}$ & $-3,5 \times 10^{-4}$ & $< 5$ & $0,1$ (CMB-S4) \\
			\bottomrule
		\end{tabular}%
}
	\end{table}
	
	\section{Vergleich mit $\Lambda$CDM}
	\label{sec:comparison}
	
	\subsection{$\chi^2$-Analyse}
	
	Vergleich der Modellanpassungen an Planck 2018-Daten:
	
	\begin{align}
		\chi^2_{\Lambda\text{CDM}} &= 1127,4 \\
		\chi^2_{T0} &= 1123,8 \\
		\Delta\chi^2 &= -3,6 \quad (1,9\sigma \text{ Verbesserung})
	\end{align}
	
	Die $\chi^2$-Werte selbst sind durch kein Prüfskript belegt.
	
	\subsection{Informationskriterien}
	
	Mit dem Akaike-Informationskriterium (AIC):
	
	\begin{equation}
		\Delta\text{AIC} = \Delta\chi^2 + 2\Delta N_{\text{params}} = -3,6 + 2 = -1,6
	\end{equation}
	
	Der negative Wert spricht in diesem Vergleich leicht für T0, trotz des zusätzlichen Parameters \textbf{[S]}.
	
	\section{Modifizierte Rekombinationsgeschichte (\texorpdfstring{$\Lambda$CDM}{LambdaCDM}-Vergleichsrahmen)}

	\begin{important}
	\textbf{Vergleichsrechnung, keine FFGFT-eigene Größe.} Das statische
	$\xi$-Universum kennt kein kosmologisches $z$ im Expansionssinn; was FFGFT
	Rotverschiebung nennt, ist fraktale Wegverlängerung, kein
	Expansions-Skalenfaktor. Die hier verwendete Rotverschiebungsvariable $z$
	ist daher eine $\Lambda$CDM-Pipeline-Größe und dient ausschließlich dem
	Zahlenvergleich mit der Standard-Rekombinationsrechnung. Sie behauptet keine
	FFGFT-Rekombinationsepoche; in FFGFT entsteht die CMB durch stationäre
	Thermalisierung in einem unendlich alten Universum (siehe $\xi$-Feld-Mechanismus
	oben), nicht durch eine Entkopplung bei irgendeinem $z$.
	\end{important}

	In diesem Vergleichsrahmen lautet die Standard-Rekombinationsbedingung:
	\begin{equation}
		z_{\text{rec}} = \text{Lösung von } x_e(z) = 0,5
	\end{equation}

	Die Elektronenfraktion entwickelt sich als:
	\begin{equation}
		x_e(z) = \frac{1}{1 + A(T) \exp[E_I/kT(z)]}
	\end{equation}

	wobei:
	\begin{align}
		T(z) &= T_0(1+z)[1 - \xi\ln(1+z)] \\
		A(T) &= \left(\frac{2\pi m_e kT}{h^2}\right)^{-3/2} 
		\frac{g_p g_e}{g_H} (1 + \xi h(T))
	\end{align}

	Mit der auf den Vergleich angewandten $\xi$-Korrektur ergibt sich
	$z_{\text{rec}} \approx 1089,5$ gegenüber
	$z_{\text{rec}}^{\Lambda\text{CDM}} = 1089,9$ in der unmodifizierten Pipeline
	-- ein Unterschied, der allein innerhalb des $\Lambda$CDM-Vergleichsrahmens
	existiert und keine FFGFT-Vorhersage einer Rekombinationsepoche ist.

	% ================== ENDE DES CMB-ABSCHNITTS ==================
	
	\section{CMB-Casimir-Verbindung und $\xi$-Feld-Verifikation}
	\label{sec:cmb_casimir}
	
	\subsection{CMB-Energiedichte und $\xi$-Längenskala}
	
	\begin{important}[These]
		\textbf{[S]} These dieses Abschnitts: Das gemessene CMB-Spektrum entspricht der strahlenden Energiedichte des $\xi$-Feld-Vakuums; das Vakuum selbst strahlt bei seiner charakteristischen Temperatur.
	\end{important}
	
	Die CMB-Energiedichte in natürlichen Einheiten:
	\begin{equation}
		\rho_{\text{CMB}} = 2,0 \times 10^{-15}\ \text{eV}^4 \quad \text{(nat. Einheiten, Dimension } [E^4] \text{)}
	\end{equation}
	
	Mit $T_{\text{CMB}} = 2,35 \times 10^{-4}$~eV ergibt sich $\rho_{\text{CMB}} = \frac{\pi^2}{15} T_{\text{CMB}}^4 = 2,0 \times 10^{-15}$~eV$^4$ und daraus $L_\xi = (\xi/\rho_{\text{CMB}})^{1/4} = 508$~eV$^{-1} = 1,00 \times 10^{-4}$~m.
	
	Die CMB-Temperatur in natürlichen Einheiten:
	\begin{equation}
		T_{\text{CMB}} = 2,35 \times 10^{-4} \quad \text{(nat. Einheiten)}
	\end{equation}
	
	Diese Energiedichte definiert eine charakteristische $\xi$-Längenskala:
	\begin{equation}
		L_\xi = \left(\frac{\xi}{\rho_{\text{CMB}}}\right)^{1/4}
	\end{equation}
	
	\begin{formula}
		Fundamentale Beziehung der CMB-Energiedichte:
		\begin{equation}
			\rho_{\text{CMB}} = \frac{\xi}{L_\xi^4} = \frac{\frac{4}{3} \times 10^{-4}}{L_\xi^4}
		\end{equation}
	\end{formula}
	
	\subsection{Casimir-CMB-Verhältnis}
	
	Der Casimir-Effekt stellt eine direkte Manifestation von Quanten-Vakuumfluktuationen dar. In natürlichen Einheiten ist die Casimir-Energiedichte zwischen zwei parallelen Platten mit Abstand $d$ (die Formel mit $240$ im Nenner beschreibt den Casimir-Druck, siehe Kraftberechnung):
	
	\begin{equation}
		|\rho_{\text{Casimir}}| = \frac{\pi^2}{720 d^4} \quad \text{(nat. Einheiten)}
	\end{equation}
	
	Bei der charakteristischen $\xi$-Längenskala $L_\xi = 10^{-4}$ m liefert das Verhältnis zwischen Casimir- und CMB-Energiedichten einen Vergleichswert:
	
	\begin{equation}
		\frac{|\rho_{\text{Casimir}}|}{\rho_{\text{CMB}}} = \frac{\pi^2}{720 \xi} = \frac{\pi^2}{720 \times \frac{4}{3} \times 10^{-4}} = \frac{\pi^2 \times 10^4}{960} \approx 102,8
	\end{equation}
	
	\subsection{Detaillierte Berechnungen in SI-Einheiten}
	
	\textbf{Casimir-Energiedichte bei Plattenabstand} $d = L_\xi = 10^{-4}$ m:
	
	\begin{align}
		|\rho_{\text{Casimir}}| &= \frac{\hbar c \pi^2}{720 d^4} \\
		&= \frac{1,055 \times 10^{-34} \times 2,998 \times 10^8 \times \pi^2}{720 \times (10^{-4})^4} \\
		&= \frac{3,12 \times 10^{-25}}{7,2 \times 10^{-14}} \\
		&= 4,3 \times 10^{-12} \text{ J/m}^3
	\end{align}
	
	\textbf{CMB-Energiedichte in SI-Einheiten:}
	\begin{equation}
		\rho_{\text{CMB}} = 4,17 \times 10^{-14} \text{ J/m}^3
	\end{equation}
	
	\textbf{Verhältnis in SI-Einheiten:}
	\begin{equation}
		\frac{|\rho_{\text{Casimir}}|}{\rho_{\text{CMB}}} = \frac{4,3 \times 10^{-12}}{4,17 \times 10^{-14}} = 104
	\end{equation}
	
	\textbf{Theoretische Vorhersage in natürlichen Einheiten:}
	\begin{align}
		\frac{|\rho_{\text{Casimir}}|}{\rho_{\text{CMB}}} &= \frac{\pi^2 / (720 L_\xi^4)}{\xi / L_\xi^4} \\
		&= \frac{\pi^2}{720 \xi} = \frac{\pi^2}{720 \times \frac{4}{3} \times 10^{-4}} \\
		&= \frac{\pi^2 \times 3 \times 10^4}{720 \times 4} = \frac{\pi^2 \times 10^4}{960} \approx 102,8
	\end{align}
	
	\textbf{Übereinstimmung:} Das in SI-Einheiten berechnete Verhältnis 104 stimmt mit dem Wert 102,8 in natürlichen Einheiten auf 1,1\% überein. Beide Werte folgen aus derselben Formel; die Nähe folgt aus der Definition von $L_\xi$ über $\rho_{\text{CMB}}$ und ist kein unabhängiger Vergleich mit einer Messung; das Verhältnis bei $L_\xi$ ist eine Identität.
	
	Die Abweichung zwischen den beiden Werten (102,8 und 104) beträgt 1,1\% und entspricht der Rundung von $L_\xi$ auf $10^{-4}$~m.
	
	\begin{important}
		Bei der charakteristischen $\xi$-Längenskala $L_\xi = 10^{-4}$ m fallen CMB-Vakuumenergiedichte und Casimir-Energiedichte nicht zusammen (Verhältnis rund $10^2$); Gleichheit liegt bei einem größeren Plattenabstand. Die Casimir-CMB-Verbindung ist eine Skalenaussage \textbf{[S]}, kein Beleg für das $\xi$-Feld.
	\end{important}
	
	\subsection{Dimensionslose $\xi$-Hierarchie und unabhängige Verifikation}
	
	\textbf{Kritische Frage: Ist dies ein Zirkelschluss?}
	
	Nach Einschätzung dieses Dokuments liegt kein Zirkelschluss vor, weil:
	
	\begin{enumerate}
		\item \textbf{Verschiedene theoretische und experimentelle Quellen:}
		\begin{itemize}
			\item $\xi$-Konstante: Rein geometrisch abgeleitet aus T0-Feldgleichungen
			\item CMB-Daten: Kosmische Mikrowellenmessungen
			\item Casimir-Messungen: Labor-Vakuumexperimente
		\end{itemize}
		
		\item \textbf{Zeitliche Abfolge der Entwicklung:}
		\begin{itemize}
			\item T0-Theorie und $\xi$-Ableitung: Rein theoretische geometrische Ableitung
			\item CMB-Vorhersage: Folgte aus der bereits etablierten $\xi$-Geometrie
			\item Casimir-Vergleich: Identität über die Definition von $L_\xi$, kein unabhängiger Pfad
		\end{itemize}
		
		\item \textbf{Mehrere Vergleichspfade:}
		\begin{itemize}
			\item Geometrische Ableitung → $\xi = \frac{4}{3} \times 10^{-4}$
			\item Higgs-Mechanismus → $\frac{\lambda_h^2 v^2}{16\pi^3 m_h^2} \approx 1{,}30 \times 10^{-4}$, rund $2{,}3\,\%$ unter $\frac{4}{3} \times 10^{-4}$ (Konsistenzprüfung, siehe Gl.~\ref{eq:xi_definition})
			\item Leptonenmassen → $\xi = \frac{4}{3} \times 10^{-4}$
			\item CMB/Casimir-Verhältnis → Identität über die Definition von $L_\xi$, kein eigener Pfad
		\end{itemize}
	\end{enumerate}
	
	\subsubsection{Detaillierte Energieskalenverhältnisse}
	
	Das dimensionslose Verhältnis zwischen CMB-Temperatur und charakteristischer Energie - detaillierte Berechnung:
	
	\begin{align}
		\frac{T_{\text{CMB}}}{E_\xi} &= \frac{2,35 \times 10^{-4}}{\frac{3}{4} \times 10^4} \\
		&= \frac{2,35 \times 10^{-4} \times 4}{3 \times 10^4} \\
		&= \frac{9,4}{3 \times 10^8} \\
		&= \frac{9,4}{3} \times 10^{-8} \\
		&= 3,13 \times 10^{-8}
	\end{align}
	
	Theoretische Vorhersage aus $\xi$-Geometrie - detaillierte Schritte:
	\begin{align}
		\xi^2 &= \left(\frac{4}{3} \times 10^{-4}\right)^2 \\
		&= \frac{16}{9} \times 10^{-8} \\
		&= 1,78 \times 10^{-8}
	\end{align}
	
	Verbesserte theoretische Vorhersage mit geometrischem Faktor:
	\begin{align}
		\frac{16}{9}\xi^2 &= \frac{16}{9} \times 1,78 \times 10^{-8} \\
		&= 1,778 \times 1,78 \times 10^{-8} \\
		&= 3,16 \times 10^{-8}
	\end{align}
	
	\textbf{Vergleich:}
	\begin{align}
		\text{Gemessen:} \quad &3,13 \times 10^{-8} \\
		\text{Theoretisch:} \quad &3,16 \times 10^{-8} \\
		\text{Übereinstimmung:} \quad &\frac{3,13}{3,16} = 0,99 = 99\% \text{ (1\% Abweichung)}
	\end{align}
	
	Übereinstimmung auf 1\% für die Beziehung
	\begin{equation}
		\boxed{\frac{T_{\text{CMB}}}{E_\xi} = \frac{16}{9}\xi^2}
	\end{equation}
	Diese Übereinstimmung ist einheitenabhängig: Die Relation ergibt eine reine Zahl, die erst in der Energieeinheit eV mit dem Messwert übereinstimmt. Sie ist daher eine Skalenwahl, keine Herleitung; die Vorwärtsherleitung der CMB-Temperatur ist offen.
	
	\subsubsection{Längenskalenverhältnisse}
	
	\begin{equation}
		\frac{\ell_{\xi}}{L_\xi} = \xi^{-1/4} = \left(\frac{3}{4}\right)^{1/4} \times 10
	\end{equation}
	
	\subsection{Konsistenz-Verifikation der T0-Theorie}
	
	\begin{important}[Konsistenzprüfung]
		Die aus der Teilchenphysik abgeleitete $\xi$-Konstante ist mit der aus der CMB gemessenen Vakuumenergiedichte verträglich; die Deutung als Vorhersage gehört zum kosmischen Sektor \textbf{[S]}.
	\end{important}
	
	Die Wege zur Längenskala im Vergleich:
	
	\begin{table}[htbp]
		\centering
		\caption{Konsistenz-Verifikation der $\xi$-Längenskala}
		\adjustbox{max width=\textwidth}{%
\begin{tabular}{p{4cm}p{4cm}p{4cm}}
			\toprule
			\textbf{Ableitung} & \textbf{Ausgangspunkt} & \textbf{Ergebnis} \\
			\midrule
			$\xi$-Geometrie (bottom-up) & $\xi = \frac{4}{3} \times 10^{-4}$ aus Teilchen & $L_\xi = 1/E_\xi = 1,33 \times 10^{-4}$ (nat. Einh.) \\
			CMB-Vakuum (top-down) & $\rho_{\text{CMB}}$ aus Messung & $L_\xi = \left(\frac{\xi}{\rho_{\text{CMB}}}\right)^{1/4} = 1,0 \times 10^{-4}$ m \\
			Casimir-Effekt & Casimir-Formel & Identität über $L_\xi$ aus $\rho_{\text{CMB}}$ \\
			\midrule
			\textbf{Vergleich} & \textbf{Zwei verschiedene Größen} & Zahlenzufall \\
			\bottomrule
		\end{tabular}%
}
	\end{table}
	Im Dokument treten zwei verschiedene Größen unter dem Namen $L_\xi$ auf. $L_\xi = 1/E_\xi = 1,33 \times 10^{-4}$ (Abschnitt zur $\xi$-Längenskala) ist eine Zahl in natürlichen Einheiten: in eV$^{-1}$ gelesen entspricht sie mit $\hbar c = 1,973 \times 10^{-7}$~eV$\cdot$m einer Länge von $2,6 \times 10^{-11}$~m, in Planck-Einheiten ($E_{\text{Planck}} = 1$, wie im Abschnitt zu den Energieskalen) $2,2 \times 10^{-39}$~m; $L_\xi = (\xi/\rho_{\text{CMB}})^{1/4} = 508$~eV$^{-1}$ entspricht $1,0 \times 10^{-4}$~m. Der Wert $10^{-4}$~m in der Zeile \enquote{bottom-up} folgt daher nicht aus $\xi$ allein, sondern aus $\rho_{\text{CMB}}$; die Übereinstimmung $1,33 \times 10^{-4} \approx 10^{-4}$ ist ein Zahlenzufall zweier verschiedener Größen.
	
	\subsection{Das $\xi$-Feld als universelles Vakuum}
	
	\begin{formula}
		Das $\xi$-Feld-Vakuum manifestiert sich in mehreren Phänomenen:
		\begin{align}
			\text{Freies Vakuum (CMB):} \quad &\rho_{\text{CMB}} = \frac{\xi}{L_\xi^4} \\
			\text{Eingeschränktes Vakuum (Casimir):} \quad &|\rho_{\text{Casimir}}| = \frac{\pi^2}{720 d^4} \\
			\text{Verhältnis bei } d = L_\xi: \quad &\frac{|\rho_{\text{Casimir}}|}{\rho_{\text{CMB}}} = \frac{\pi^2 \times 10^4}{960}
		\end{align}
	\end{formula}
	
	\begin{important}
		Alle $\xi$-Beziehungen bestehen aus exakten mathematischen Verhältnissen:
		\begin{itemize}
			\item Brüche: $\frac{4}{3}$, $\frac{16}{9}$, $\frac{3}{4}$
			\item Zehnerpotenzen: $10^{-4}$, $10^4$
			\item Mathematische Konstanten: $\pi^2$
		\end{itemize}
		Es treten keine frei angepassten Dezimalzahlen auf; die Verhältnisse folgen aus der $\xi$-Geometrie.
	\end{important}
	
	\section{Casimir-Effekt und $\xi$-Feld-Verbindung}
	
	\subsection{Modifizierte Casimir-Formel in der T0-Theorie}
	
	Die T0-Theorie bietet eine Deutung des Casimir-Effekts durch das $\xi$-Feld:
	
	\begin{equation}
		|\rho_{\text{Casimir}}(d)| = \frac{\pi^2}{720 \xi} \rho_{\text{CMB}} \left(\frac{L_\xi}{d}\right)^4
	\end{equation}
	
	Einsetzen von $\rho_{\text{CMB}} = \xi/L_\xi^4$ ergibt die Standardformel:
	\begin{equation}
		|\rho_{\text{Casimir}}| = \frac{\pi^2}{720 d^4}
	\end{equation}
	
	Das ist mit der Deutung verträglich, dass der Casimir-Effekt und die CMB verschiedene Manifestationen desselben $\xi$-Feld-Vakuums sind \textbf{[S]}.
	
	\section{Strukturbildung im statischen $\xi$-Universum}
	
	\subsection{Kontinuierliche Strukturentwicklung}
	
	Im statischen T0-Universum \textbf{[S]} findet Strukturbildung kontinuierlich ohne Urknall-Einschränkungen statt:
	
	\begin{equation}
		\frac{d\rho}{dt} = -\nabla \cdot (\rho \mathbf{v}) + S_\xi(\rho, T, \xi)
	\end{equation}
	
	wobei $S_\xi$ der $\xi$-Feld-Quellterm für kontinuierliche Materie/Energie-Transformation ist.
	
	\subsection{$\xi$-unterstützte kontinuierliche Schöpfung}
	
	In diesem Bild \textbf{[S]} ermöglicht das $\xi$-Feld kontinuierliche Materie/Energie-Transformation:
	
	\begin{align}
		\text{Quantenvakuum} &\xrightarrow{\xi} \text{Virtuelle Teilchen} \\
		\text{Virtuelle Teilchen} &\xrightarrow{\xi^2} \text{Reale Teilchen} \\
		\text{Reale Teilchen} &\xrightarrow{\xi^3} \text{Atomkerne} \\
		\text{Atomkerne} &\xrightarrow{\text{Zeit}} \text{Sterne, Galaxien}
	\end{align}
	
	Die Energiebilanz wird aufrechterhalten durch:
	\begin{equation}
		\rho_{\text{total}} = \rho_{\text{Materie}} + \rho_{\xi\text{-Feld}} = \text{konstant}
	\end{equation}
	
	\begin{important}
		\textbf{[S]} In diesem Bild bleibt die Energie durch kontinuierliche Transformation zwischen Materie und $\xi$-Feld-Energie erhalten, was eine ewige Existenz ohne Anfang oder Ende zulässt.
	\end{important}
	
	\section{Einheitenanalyse der $\xi$-basierten Casimir-Formel}
	
	Diese Analyse untersucht die Einheitenkonsistenz der modifizierten Casimir-Formel innerhalb der T0-Theorie, die die dimensionslose Konstante $\xi$ und die kosmische Mikrowellen-Hintergrund-(CMB)-Energiedichte $\rho_{\text{CMB}}$ einführt. Das Ziel ist, die Konsistenz mit der Standard-Casimir-Formel zu verifizieren und die physikalische Bedeutung der neuen Parameter $\xi$ und $L_\xi$ zu klären. Die Analyse wird in SI-Einheiten durchgeführt, wobei jede Formel auf dimensionale Korrektheit geprüft wird.
	
	\subsection{Standard-Casimir-Formel}
	Die Standard-Casimir-Formel beschreibt die Energiedichte des Casimir-Effekts zwischen zwei parallelen, perfekt leitenden Platten im Vakuum:
	\begin{equation}
		|\rho_{\text{Casimir}}| = \frac{\pi^2 \hbar c}{720 d^4}
	\end{equation}
	Hier ist $\hbar$ die reduzierte Planck-Konstante, $c$ die Lichtgeschwindigkeit und $d$ der Abstand zwischen den Platten. Die Einheitenprüfung ergibt:
	\begin{equation}
		\frac{[\hbar] \cdot [c]}{[d^4]} = \frac{(\text{J} \cdot \text{s}) \cdot (\text{m}/\text{s})}{\text{m}^4} = \frac{\text{J} \cdot \text{m}}{\text{m}^4} = \frac{\text{J}}{\text{m}^3}
	\end{equation}
	Dies entspricht der Einheit der Energiedichte und bestätigt die Korrektheit der Formel.
	
	\textbf{Formelerklärung:} Der Casimir-Effekt entsteht aus Quantenfluktuationen des elektromagnetischen Feldes im Vakuum. Nur bestimmte Wellenlängen passen zwischen die Platten, was zu einer messbaren Energiedichte führt, die mit $d^{-4}$ skaliert. Die Konstante $\pi^2/720$ ergibt sich aus der Summierung über alle erlaubten Moden.
	
	\subsection{Definition von $\xi$ und CMB-Energiedichte}
	Die T0-Theorie führt die dimensionslose Konstante $\xi$ ein, definiert als:
	\begin{equation}
		\xi = \frac{4}{3} \times 10^{-4}
	\end{equation}
	Diese Konstante ist dimensionslos, bestätigt durch $[\xi] = [1]$. Die CMB-Energiedichte ist in natürlichen Einheiten definiert als:
	\begin{equation}
		\rho_{\text{CMB}} = \frac{\xi}{L_\xi^4}
	\end{equation}
	mit der charakteristischen Längenskala $L_\xi = 10^{-4}$ m. In SI-Einheiten ist die CMB-Energiedichte:
	\begin{equation}
		\rho_{\text{CMB}} = 4,17 \times 10^{-14} \text{ J}/\text{m}^3
	\end{equation}
	
	\textbf{Formelerklärung:} Die CMB-Energiedichte repräsentiert die Energie der kosmischen Mikrowellen-Hintergrundstrahlung. In der T0-Theorie wird sie durch $\xi$ und $L_\xi$ skaliert, wobei $L_\xi$ eine fundamentale Längenskala ist, die möglicherweise mit kosmischen Phänomenen verknüpft ist. Die Einheitenanalyse zeigt:
	\begin{equation}
		[\rho_{\text{CMB}}] = \frac{[\xi]}{[L_\xi^4]} = \frac{1}{\text{m}^4} = \text{E}^4 \text{ (in natürlichen Einheiten)}
	\end{equation}
	In SI-Einheiten ergibt dies J/m$^3$, was konsistent ist.
	
	\subsection{Konversion der $\xi$-Beziehung zu SI-Einheiten}
	Die T0-Theorie postuliert eine fundamentale Beziehung:
	\begin{equation}
		\hbar c \stackrel{!}{=} \frac{\rho_{\text{CMB}} L_\xi^4}{\xi}
	\end{equation}
	Die Einheitenanalyse bestätigt:
	\begin{equation}
		\frac{[\rho_{\text{CMB}}] \cdot [L_\xi^4]}{[\xi]} = \left( \frac{\text{J}}{\text{m}^3} \right) \cdot \text{m}^4 \cdot 1 = \text{J} \cdot \text{m}
	\end{equation}
	Dies entspricht der Einheit von $\hbar c$. Numerisch erhalten wir:
	\begin{equation}
		\frac{\left( 4,17 \times 10^{-14} \right) \cdot \left( 10^{-4} \right)^4}{\frac{4}{3} \times 10^{-4}} = 3,13 \times 10^{-26} \text{ J} \cdot \text{m}
	\end{equation}
	
	Die Beziehung folgt aus der Definition $\rho_{\text{CMB}} = \xi\hbar c/L_\xi^4$. Verglichen mit $\hbar c = 3,16 \times 10^{-26}$ J·m ist das das $0,99$-fache; die Übereinstimmung auf etwa $1\,\%$ folgt aus der Festlegung von $L_\xi$ über $\rho_{\text{CMB}}$.
	
	\textbf{Formelerklärung:} Diese Beziehung überbrückt Quantenmechanik ($\hbar c$) mit kosmischen Skalen ($\rho_{\text{CMB}}$, $L_\xi$). Die dimensionslose Konstante $\xi$ fungiert als Skalierungsfaktor, der die CMB-Energiedichte mit der fundamentalen Längenskala $L_\xi$ verknüpft.
	
	\subsection{Modifizierte Casimir-Formel}
	Die modifizierte Casimir-Formel ist:
	\begin{equation}
		|\rho_{\text{Casimir}}(d)| = \frac{\pi^2}{720 \xi} \rho_{\text{CMB}} \left( \frac{L_\xi}{d} \right)^4
	\end{equation}
	Die Einheitenanalyse ergibt:
	\begin{equation}
		\frac{[\rho_{\text{CMB}}] \cdot [L_\xi^4]}{[\xi] \cdot [d^4]} = \frac{\left( \frac{\text{J}}{\text{m}^3} \right) \cdot \text{m}^4}{1 \cdot \text{m}^4} = \frac{\text{J}}{\text{m}^3}
	\end{equation}
	Dies bestätigt die Einheit der Energiedichte. Einsetzen von $\rho_{\text{CMB}} = \xi \hbar c / L_\xi^4$ ergibt die Standard-Casimir-Formel:
	\begin{equation}
		|\rho_{\text{Casimir}}| = \frac{\pi^2}{720} \frac{\xi \hbar c}{L_\xi^4} \cdot \frac{L_\xi^4}{d^4} = \frac{\pi^2 \hbar c}{720 d^4}
	\end{equation}
	
	\textbf{Formelerklärung:} Die modifizierte Formel beinhaltet $\xi$ und $\rho_{\text{CMB}}$, was den Casimir-Effekt mit kosmischen Parametern verknüpft. Ihre Konsistenz mit der Standardformel zeigt, dass die T0-Theorie eine alternative Darstellung des Effekts bietet.
	
	\subsection{Kraftberechnung}
	Die Kraft pro Fläche wird aus der Energiedichte abgeleitet:
	\begin{equation}
		\frac{F}{A} = -\frac{\partial}{\partial d} \left( |\rho_{\text{Casimir}}| \cdot d \right) = \frac{\pi^2}{240 \xi} \rho_{\text{CMB}} \left( \frac{L_\xi}{d} \right)^4
	\end{equation}
	
	Mit $\rho_{\text{CMB}} = \xi\hbar c/L_\xi^4$ ergibt sich daraus die Literaturformel des Casimir-Drucks $\pi^2\hbar c/(240\,d^4)$.
	Die Einheitenanalyse zeigt:
	\begin{equation}
		\frac{[\rho_{\text{CMB}}] \cdot [L_\xi^4]}{[\xi] \cdot [d^4]} = \frac{\left( \frac{\text{J}}{\text{m}^3} \right) \cdot \text{m}^4}{1 \cdot \text{m}^4} = \frac{\text{J}}{\text{m}^3} = \frac{\text{N}}{\text{m}^2}
	\end{equation}
	Dies entspricht der Einheit des Drucks und bestätigt die Korrektheit.
	
	\textbf{Formelerklärung:} Die Kraft pro Fläche repräsentiert die messbare Casimir-Kraft, die aus der Änderung der Energiedichte mit dem Plattenabstand entsteht. Die T0-Theorie skaliert diese Kraft mit $\xi$ und $\rho_{\text{CMB}}$, was eine kosmische Interpretation ermöglicht.
	
	\subsection{Kritische Bewertung}
	Die T0-Theorie zeigt Einheitenkonsistenz; $\rho_{\text{CMB}} L_\xi^4/\xi$ trifft $\hbar c$ auf etwa $1\,\%$, was aus der Festlegung von $L_\xi$ folgt (ein Faktor 16/9 tritt dabei nicht auf). Sie verknüpft den Casimir-Effekt mit kosmischer Vakuumenergie über $\xi$ und $L_\xi$, wobei $L_\xi = 10^{-4}$ m als fundamentale Längenskala fungiert. Dies eröffnet neue physikalische Interpretationen, die den Casimir-Effekt mit kosmologischen Phänomenen verbinden.
	
	\section{Dimensionslose $\xi$-Hierarchie}
	
	\subsection{Vollständige Tabelle dimensionsloser Verhältnisse}
	
	Alle $\xi$-Beziehungen reduzieren sich auf exakte mathematische Verhältnisse:
	
	\begin{table}[htbp]
		\centering
		\caption{Dimensionslose $\xi$-Verhältnisse in der T0-Theorie}
		\adjustbox{max width=\textwidth}{%
\begin{tabular}{lcc}
			\toprule
			\textbf{Verhältnis} & \textbf{Ausdruck} & \textbf{Wert} \\
			\midrule
			Temperaturverhältnis & $\frac{T_{\text{CMB}}}{E_\xi}$ & $3,13 \times 10^{-8}$ \\
			Relation (einheitenabhängig) & $\frac{16}{9}\xi^2$ & $3,16 \times 10^{-8}$ \\
			Längenverhältnis & $\frac{\ell_{\xi}}{L_\xi}$ & $\xi^{-1/4}$ \\
			Casimir-CMB & $\frac{|\rho_{\text{Casimir}}|}{\rho_{\text{CMB}}}$ & $\frac{\pi^2 \times 10^4}{960}$ \\
			Gravitationskopplung & $\alpha_G$ & $\xi^2 = 1,78 \times 10^{-8}$ \\
			Schwache Kopplung & $\alpha_W$ & $\xi^{1/2} = 1,15 \times 10^{-2}$ \\
			Starke Kopplung & $\alpha_S$ & $\xi^{-1/3} = 19,57$ \\
			\bottomrule
		\end{tabular}%
}
	\end{table}
	
	\begin{important}
		Alle $\xi$-Beziehungen bestehen aus exakten mathematischen Verhältnissen:
		\begin{itemize}
			\item Brüche: $\frac{4}{3}$, $\frac{3}{4}$, $\frac{16}{9}$
			\item Zehnerpotenzen: $10^{-4}$, $10^3$, $10^4$
			\item Mathematische Konstanten: $\pi^2$
		\end{itemize}
		Es treten keine frei angepassten Dezimalzahlen auf; die Verhältnisse folgen aus der $\xi$-Geometrie.
	\end{important}
	
	\subsection{Parameterreduktion}
	
	\begin{important}[Parameterzahl]
		Die T0-Theorie arbeitet mit dem einzigen Parameter $\xi$ (plus den im Korpus genannten ganzzahligen Setzungen; SI-/SM-Anker wie $m_e$ dienen der Umrechnung):
		\begin{itemize}
			\item Standardmodell der Teilchenphysik: 19+ Parameter
			\item $\Lambda$CDM-Kosmologie: 6 Parameter
			\item T0-Theorie: 1 Parameter ($\xi$)
		\end{itemize}
		Das entspricht einer Reduktion der Parameterzahl um 96\%; der Vergleich mit $\Lambda$CDM gehört zum kosmischen Sektor und ist spekulativ \textbf{[S]} (P39).
	\end{important}
	
	\section{Einheitenanalyse und dimensionale Konsistenz}
	
	\subsection{Verifikation des Rahmenwerks natürlicher Einheiten}
	
	Alle T0-Theorie-Gleichungen sind in natürlichen Einheiten dimensional konsistent:
	
	\begin{table}[h]
		\centering
		\adjustbox{max width=\textwidth}{%
\begin{tabular}{l l l l}
			\toprule
			Größ{}e & Natürliche Einheiten & Dimension & Verifikation \\
			\midrule
			$\xi$ & dimensionslos & $[1]$ & erfüllt \\
			$E_\xi$ & 7500 & $[E]$ & erfüllt \\
			$L_\xi$ & $1,33 \times 10^{-4}$ & $[E^{-1}]$ & erfüllt \\
			$T_\xi$ & 7500 & $[E]$ & erfüllt \\
			$G_{\text{nat}}$ & $2,61 \times 10^{-70}$ & $[E^{-2}]$ & erfüllt \\
			\bottomrule
		\end{tabular}%
}
		\caption{Dimensionale Konsistenz in natürlichen Einheiten}
	\end{table}
	
	\subsection{Energieskalen-Hierarchien}
	
	Mit der Planck-Energie als Bezug gilt:
	
	\begin{align}
		E_{\text{Planck}} &= 1 \quad \text{(per Definition in natürlichen Einheiten)} \\
		E_\xi &= \frac{1}{\xi} = 7500
	\end{align}
	Die einfachen Potenzen $\xi^{1/2}\,E_{\text{Planck}} \approx 0,0115\,E_P = 1,4 \times 10^{17}$~GeV und $\xi^{1/3}\,E_{\text{Planck}} \approx 0,0511\,E_P = 6,2 \times 10^{17}$~GeV (mit $E_{\text{Planck}} = 1,22 \times 10^{19}$~GeV) liefern dagegen nicht die schwache und die QCD-Skala: Gemessen sind $v/E_P = 246/1,22 \times 10^{19} = 2,0 \times 10^{-17}$ (im Korpus: $v/E_P = \xi^4/(5\pi)$, Dok.~149) und $\Lambda_{\text{QCD}}/E_P \approx 0,2/1,22 \times 10^{19} = 1,6 \times 10^{-20}$; die Potenzen $\xi^{1/2}$ und $\xi^{1/3}$ verfehlen diese Skalen um 15 bzw.\ 18 Größenordnungen.
	
	\subsection{Zusätzliche experimentelle Vorhersagen}
	
	\textbf{Vorhersage 1: Elektromagnetische Resonanz bei charakteristischer $\xi$-Frequenz}
	\begin{itemize}
		\item Maximale $\xi$-Feld-Photon-Kopplung bei $\nu = E_\xi = 7500$ (nat. Einheiten)
		\item Anomalien in elektromagnetischer Ausbreitung bei dieser Frequenz
		\item Spektrale Besonderheiten im entsprechenden Frequenzbereich
	\end{itemize}
	
	\textbf{Vorhersage 2: Casimir-Kraft-Anomalien bei charakteristischer $\xi$-Längenskala}
	\begin{itemize}
		\item Standard-Casimir-Gesetz: $F \propto d^{-4}$
		\item $\xi$-Feld-Modifikationen bei $d \approx L_\xi = 10^{-4}$ m
		\item Messbare Abweichungen durch $\xi$-Vakuum-Kopplung
	\end{itemize}
	
	\textbf{Vorhersage 3: Modifizierte Vakuumfluktuationen}
	\begin{itemize}
		\item Vakuumenergiedichte-Variationen bei Skala $L_\xi$
		\item Korrelation zwischen Casimir- und CMB-Messungen
		\item Testbar in Präzisions-Laborexperimenten
	\end{itemize}
	
	\section{Das statische Universums-Paradigma}
	
	\subsection{Fundamentale Eigenschaften des T0-Universums}
	
	\begin{important}[Kosmologisches Bild]
		\textbf{[S]} Das T0-Universum unterscheidet sich in folgenden Punkten von der Expansionskosmologie (spekulativ, P39):
		\begin{itemize}
			\item Das Universum expandiert nicht
			\item Das Universum hat ewig existiert
			\item Das Universum hat keinen Anfang (kein Urknall)
			\item Das Universum befindet sich im thermodynamischen Gleichgewicht
			\item Kosmische Phänomene werden auf $\xi$-Feld-Dynamik zurückgeführt
		\end{itemize}
	\end{important}
	
	\subsection{$r_0$-Definition aus $\xi$}
	
	Die fundamentale Längenskala $r_0$ ist definiert durch:
	\begin{align}
		r_0 &= \xi \cdot l_P = \frac{4}{3} \times 10^{-4} \times 1,616 \times 10^{-35}\,\text{m} \\
		&= 2,15 \times 10^{-39}\,\text{m}
	\end{align}
	
	In natürlichen Einheiten mit $l_P = 1$:
	\begin{equation}
		r_0 = \xi = \frac{4}{3} \times 10^{-4}
	\end{equation}
	
	\section{Die zentrale These: Das Vakuum ist das $\xi$-Feld}
	
	\begin{formula}
		Die universelle $\xi$-Konstante führt auf folgende zusammenhängende Beziehungen:
		\begin{align}
			\xi &= \frac{4}{3} \times 10^{-4} \quad \text{(aus Geometrie)} \\
			G &= \frac{\xi^2}{4m} \quad \text{(Gravitation berechenbar)} \\
			T_{\text{CMB}} &= \frac{16}{9} \xi^2 \times E_\xi \quad \text{(einheitenabhängig auf ca.\ 1\,\% getroffen, s.\ oben)} \\
			\frac{|\rho_{\text{Casimir}}|}{\rho_{\text{CMB}}} &= \frac{\pi^2 \times 10^4}{960} \quad \text{(Casimir-Verbindung)}
		\end{align}
	\end{formula}
	
	\subsection{Das Vakuum ist das $\xi$-Feld}
	
	\begin{important}
		Zentrale These der T0-Theorie:
		\begin{itemize}
			\item Das Vakuum ist identisch mit dem $\xi$-Feld
			\item Die CMB ist Strahlung dieses Vakuums bei charakteristischer Temperatur \textbf{[S]}
			\item Die Casimir-Kraft entsteht aus geometrischer Einschränkung desselben Vakuums
			\item Gravitation folgt aus $\xi$-Geometrie
			\item Die fundamentalen Kräfte werden auf $\xi$-Feld-Manifestationen zurückgeführt
		\end{itemize}
	\end{important}
	
	\subsection{Mathematische Struktur}
	
	Die T0-Theorie formuliert:
	\begin{enumerate}
		\item \textbf{Universelle $\xi$-Skalierung}: Die behandelten Phänomene werden aus $\xi = \frac{4}{3} \times 10^{-4}$ hergeleitet
		\item \textbf{Statisches Paradigma} \textbf{[S]}: Kein Urknall, keine Expansion, ewige Existenz
		\item \textbf{Zeit-Energie-Konsistenz}: $T\cdot E = 1$ mit unterer Zeitschranke $T_0 = \xi\,t_P$
		\item \textbf{Dimensionale Konsistenz}: Vollständig formuliert in natürlichen Einheiten
		\item \textbf{Einheiten-unabhängige Physik}: Exakte mathematische Verhältnisse
	\end{enumerate}
	
	\section{Schlussfolgerungen}
	
	Die T0-Analyse der Temperatureinheiten in natürlichen Einheiten mit CMB-Berechnungen ergibt:
	
	\begin{enumerate}
		\item \textbf{Universelle $\xi$-Skalierung}: Die behandelten Temperatur- und Energieskalen folgen aus der geometrischen Konstante $\xi = \frac{4}{3} \times 10^{-4}$.
		
		\item \textbf{CMB ohne Inflation} \textbf{[S]}: Die Theorie deutet die CMB ohne Inflation und ohne eine Entkopplung bei irgendeinem kosmologischen $z$ (stationäre $\xi$-Feld-Thermalisierung in einem unendlich alten Universum) und leitet primordiale Störungen aus T-Feld-Quantenfluktuationen ab.
		
		\item \textbf{Statisches Universums-Paradigma} \textbf{[S]}: In diesem Bild ist das Universum ewig und statisch und mit der Quantenmechanik ohne Paradoxe vereinbar.
		
		\item \textbf{Zeit-Energie-Konsistenz} \textbf{[S]}: Aus $T\cdot E = 1$ mit der unteren Zeitschranke $T_0 = \xi\,t_P$ und dem randlosen Torus folgt, dass es keinen singulären Anfang gibt (Dok.~306); das statische Universum benötigt keinen Urknall.
		
		\item \textbf{Dimensionale Konsistenz}: Durchgehend konsistente Dimensionen in natürlichen Einheiten, mit $\xi$ als einzigem Parameter (plus den genannten geometrischen Setzungen).
		
		\item \textbf{Einheiten-unabhängige Physik}: Alle Beziehungen bestehen aus exakten mathematischen Verhältnissen, die aus fundamentaler Geometrie abgeleitet sind.
		
		\item \textbf{Testbare Vorhersagen}: Spezifische, messbare Abweichungen vom $\Lambda$CDM, die mit Experimenten der nächsten Generation getestet werden können.
	\end{enumerate}
	
	\begin{important}[Fazit]
		Die T0-Theorie schlägt eine in natürlichen Einheiten formulierte Alternative zur expansionsbasierten Kosmologie vor \textbf{[S]} und beschreibt Temperaturphänomene von der Teilchenphysik bis zum Kosmos mit einer einzigen Konstante, die aus Geometrie abgeleitet wird. Die CMB-Berechnungen zeigen, dass kosmologische Beobachtungen innerhalb dieses Rahmens beschrieben werden können; der kosmologische Teil bleibt spekulativ (P39).
	\end{important}
	

	
	\begin{thebibliography}{20}
		\bibitem{T0Theory}
		Johann Pascher.
		\textit{Das T0-Modell (Planck-referenziert): Eine Neuformulierung der Physik}.
		GitHub Repository, 2024.
		\url{https://jpascher.github.io/T0-Time-Mass-Duality/2/pdf}
		
		\bibitem{FineStructure}
		Johann Pascher.
		\textit{Die Feinstrukturkonstante: Verschiedene Darstellungen und Beziehungen}.
		Erklärt die kritische Unterscheidung zwischen $\alpha_{\text{EM}} = 1/137$ (SI) und $\alpha_{\text{EM}} = 1$ (natürliche Einheiten).
		2025.
		
		\bibitem{planck2020}
		Planck Collaboration (2020). 
		\textit{Planck 2018 Ergebnisse. VI. Kosmologische Parameter}. 
		Astronomy \& Astrophysics, 641, A6. 
		\url{https://doi.org/10.1051/0004-6361/201833910}
		
		\bibitem{codata2018}
		CODATA (2018). 
		\textit{Die 2018 CODATA empfohlenen Werte der fundamentalen physikalischen Konstanten}. 
		National Institute of Standards and Technology. 
		\url{https://physics.nist.gov/cuu/Constants/}
		
		\bibitem{casimir1948}
		Casimir, H. B. G. (1948). 
		\textit{Über die Anziehung zwischen zwei perfekt leitenden Platten}. 
		Proceedings of the Royal Netherlands Academy of Arts and Sciences, 51(7), 793--795.
		
		\bibitem{muon_g2_2021}
		Myon g-2 Kollaboration (2021). 
		\textit{Messung des positiven Myon anomalen magnetischen Moments auf 0,46 ppm}. 
		Physical Review Letters, 126(14), 141801. 
		\url{https://doi.org/10.1103/PhysRevLett.126.141801}
		
		\bibitem{riess2022}
		Riess, A. G., et al. (2022). 
		\textit{Eine umfassende Messung des lokalen Wertes der Hubble-Konstante mit 1 km s$^{-1}$ Mpc$^{-1}$ Unsicherheit vom Hubble-Weltraumteleskop und dem SH0ES-Team}. 
		The Astrophysical Journal Letters, 934(1), L7. 
		\url{https://doi.org/10.3847/2041-8213/ac5c5b}
		
		\bibitem{jwst_early}
		Naidu, R. P., et al. (2022). 
		\textit{Zwei bemerkenswert leuchtende Galaxienkandidaten bei z $\approx$ 11--13 enthüllt durch JWST}. 
		The Astrophysical Journal Letters, 940(1), L14. 
		\url{https://doi.org/10.3847/2041-8213/ac9b22}
		
		\bibitem{cobe1992}
		COBE Kollaboration (1992). 
		\textit{Struktur in den COBE Differential-Mikrowellen-Radiometer Erstkarten}. 
		The Astrophysical Journal Letters, 396, L1--L5. 
		\url{https://doi.org/10.1086/186504}
	\end{thebibliography}
